\scanpage{79}
% archive foliation: 78
\[
(20)\qquad \alpha:\underline{\omega}_J\simeq\mu_4^*\qquad\text{i.e. d'un élément de }\underline{\omega}_J\wedge\mu_4^*.
\]

Ces définitions posées, on voit donc que la catégorie\inserted{-groupoïde} des courbes elliptiques\inserted{/$k$}, épinglées en Legendre \deleted{\unsure{sera}}\inserted{est} équivalente à la catégorie-groupoïde des \deleted{triplets couples}\inserted{triplets} $(J,t,\alpha)$ \deleted{de} formés
\[
(21)\quad\left\{\begin{array}{ll}
a)&\text{d'un ens }J\text{ à trois éléments (d'où }\Sigma_J,D_J,\Sigma_J^*\text{)}\\[4pt]
b)&\text{d'un }t\in\Sigma_J^*(k)\\[4pt]
c)&\text{d'un isom }\alpha:\underline{\omega}_J\simeq\mu_4^*\deleted{\unsure{k}}\quad(=\mu_4^*(k))
\end{array}\right.
\]

Un épinglage de Legendre-Jacobi d'échelon $2$ \inserted{d'une courbe elliptique $E$} sera la donnée d'un couple d'épinglages, l'un de Legendre, l'autre de Jacobi d'échelon $2$ i.e. d'une bijection $J\simeq\{0,1,\infty\}$, permettant d'identifier $\Sigma_J$ à la droite projective type $\mathbb{P}^1_{\unsure{k}}$. On a alors dans $\underline{\omega}_J\simeq\underline{\omega}_{\{0,1,\infty\}}$ un élément privilégié $\omega^+$, et la donnée c) \deleted{\unsure{du choix}} de l'épinglage de Legendre équivaut\inserted{donc} à celle d'une racine carrée \deleted{\unsure{choisie}} de $-1$, soit $i$ --- c'est la donnée ``arithmétique''. \deleted{Par ailleurs, la catég} La catégorie-groupoïde des courbes elliptiques sur $k$, munies d'un épinglage de Legendre-Jacobi, \deleted{équivaut à}\inserted{est rigide (pas d'auto. non identiques)} \deleted{celle des \unsure{objets} à isom. près \unsure{ce}} et elle
